\chapter*{Abstract}
\label{abstract}

Cloud computing has become a fundamental infrastructure for modern software development by providing scalable automated on-demand services. Furthermore, hybrid cloud has gained importance thanks to its characteristic including the flexibility of public-cloud services with the security of private infrastructure. However, operating these heterogeneous environments introduces challenges in maintaining consistent configurations, reconciling infrastructure changes, and enforcing security policies. Moreover, existing practices often separate system administration, software development, and security assessment into different tools and workflows, which can delay feedback and produce distinct findings. Although existing research has improved individual cloud capabilities through many approaches, specifically DevOps, and security assessment tools, less research has been given to combining these inventions into an operational process. This research aims to design and develop a SysSecOps operational process that merges infrastructure management, continuous security, and software delivery within a hybrid cloud environment. The proposed process evaluates infrastructure, aggregates and unite findings from multiple assessment sources, calculates a risk score, applies a pass, review, or reject decision for deployment,and monitoring the applications. In addition, a labelled dataset containing vulnerable and reference secure code was constructed to train and evaluate a code-risk model to distinguish vulnerabilities from safe samples. The results indicate that the proposed process has the potential in scalable cloud operational tasks, while the risk-scoring model can provide useful suggestions for identifying and prioritising vulnerabilities. The research provides a foundation for more consistent, traceable, and risk-aware operations for hybrid cloud infrastructure.


\endinput
% Bản Tiếng Việt
Điện toán đám mây đã trở thành nền tảng cho sự phát triển phần mềm hiện đại nhờ khả năng cung cấp các dịch vụ tự động, có thể mở rộng theo nhu cầu. Bên cạnh đó, mô hình đám mây lai ngày càng trở nên quan trọng nhờ kết hợp được tính linh hoạt của các dịch vụ đám mây công cộng với mức độ bảo mật của hạ tầng riêng. Tuy nhiên, vận hành các môi trường này đặt ra nhiều thách thức trong việc duy trì cấu hình nhất quán, đồng bộ thay đổi hạ tầng và thực thi chính sách bảo mật. Hơn nữa, các phương pháp hiện hành thường tách hoạt động quản trị hệ thống, phát triển phần mềm và đánh giá bảo mật thành các công cụ và quy trình làm việc chuyên biệt. Điều này có thể làm chậm quá trình phản hồi và tạo ra các kết quả đánh giá rời rạc, thiếu nhất quán. Mặc dù các nghiên cứu hiện nay đã cải thiện từng năng lực riêng lẻ của hệ thống đám mây thông qua nhiều phương pháp, đặc biệt là DevOps, và các công cụ đánh giá bảo mật, song vẫn chưa có nhiều nghiên cứu tập trung vào việc kết hợp những tiến bộ này thành một quy trình vận hành thống nhất. Nghiên cứu này nhằm thiết kế và phát triển một quy trình vận hành SysSecOps, tích hợp quản lý hạ tầng, bảo mật liên tục và phân phối phần mềm trong môi trường đám mây lai. Quy trình được đề xuất thực hiện đánh giá hạ tầng, tổng hợp và hợp nhất các phát hiện từ nhiều nguồn đánh giá, tính toán điểm rủi ro, đưa ra quyết định cho phép triển khai, yêu cầu xem xét hoặc từ chối triển khai, đồng thời giám sát các ứng dụng sau khi triển khai. Hai kịch bản thực tế đã được xây dựng nhằm giải quyết các thách thức nêu trên và được đề xuất trong một môi trường vận hành thực tế. Bên cạnh đó, một tập dữ liệu có gán nhãn, bao gồm mã nguồn chứa lỗ hổng và mã nguồn an toàn đã được xây dựng để huấn luyện và đánh giá mô hình chấm điểm rủi ro trong việc phân biệt các mã độc với các mã nguồn an toàn. Kết quả nghiên cứu cho thấy quy trình được đề xuất có tiềm năng hỗ trợ các tác vụ vận hành đám mây có khả năng mở rộng, trong khi mô hình chấm điểm rủi ro có thể cung cấp các đề xuất hữu ích để xác định và ưu tiên xử lý các lỗ hổng bảo mật. Nghiên cứu này tạo nền tảng cho việc xây dựng các hoạt động vận hành hạ tầng đám mây lai nhất quán hơn, có khả năng truy vết và dựa trên nhận thức về rủi ro.


Giới hạn 200-300 từ, sử dụng văn phong học thuật, ngắn gọn, khách quan, viết liền mạch, không gạch đầu dòng.

Hướng dẫn trình bày tóm tắt:

\begin{enumerate}
    \item Giới thiệu vấn đề/ngữ cảnh của vấn đề (khoảng 1 - 3 câu): Trình bày bối cảnh và vấn đề thực tế đang tồn tại/ khoảng trống nghiên cứu, tầm quan trọng của vấn đề nghiên cứu, lý do cần thực hiện đề tài.
    \item Mục tiêu đề tài (1 câu): nêu cụ thể đề tài giải quyết/phát triển điều gì.
    \item Phương pháp và giải pháp (1 - 3 câu): Trình bày phương pháp, công nghệ, công cụ, mô hình, giải pháp được đề xuất trong đề tài. Nêu những cải tiến nổi bật (nếu có). Nêu phương pháp thử nghiệm, đánh giá, bộ dữ liệu được sử dụng để đánh giá (nếu có)
    \item Kết quả chính (1 - 2 câu): Trình bày ngắn gọn các kết quả nổi bật, có thể đưa số liệu cụ thể (độ chính xác, tốc độ xử lý, dung lượng dữ liệu, \ldots)
    \item Kết luận (1 - 2 câu): Nêu ý nghĩa khoa học và thực tiễn (nếu có) của phương pháp/giải pháp hay kết quả của đề tài, khả năng áp dụng vào thực tế/đóng góp cho doanh nghiệp/người dùng/nghiên cứu sau này, \ldots
\end{enumerate}